% Options for packages loaded elsewhere
% Options for packages loaded elsewhere
\PassOptionsToPackage{unicode}{hyperref}
\PassOptionsToPackage{hyphens}{url}
\PassOptionsToPackage{dvipsnames,svgnames,x11names}{xcolor}
%
\documentclass[
  letterpaper,
  DIV=11,
  numbers=noendperiod]{scrartcl}
\usepackage{xcolor}
\usepackage{amsmath,amssymb}
\setcounter{secnumdepth}{-\maxdimen} % remove section numbering

\usepackage{fontspec}
\usepackage{unicode-math}

% Fuente principal
\setmainfont{Times New Roman}

% Fuente matemática compatible
\setmathfont{TeX Gyre Termes Math}

% Hacer que KOMA-Script también use la fuente romana
% en títulos y encabezados
\setkomafont{disposition}{\rmfamily\bfseries}
\setkomafont{title}{\rmfamily\bfseries}
\setkomafont{author}{\rmfamily}
\setkomafont{date}{\rmfamily}

% Use upquote if available, for straight quotes in verbatim environments
\IfFileExists{upquote.sty}{\usepackage{upquote}}{}
\IfFileExists{microtype.sty}{% use microtype if available
  \usepackage[]{microtype}
  \UseMicrotypeSet[protrusion]{basicmath} % disable protrusion for tt fonts
}{}
\makeatletter
\@ifundefined{KOMAClassName}{% if non-KOMA class
  \IfFileExists{parskip.sty}{%
    \usepackage{parskip}
  }{% else
    \setlength{\parindent}{0pt}
    \setlength{\parskip}{6pt plus 2pt minus 1pt}}
}{% if KOMA class
  \KOMAoptions{parskip=half}}
\makeatother
% Make \paragraph and \subparagraph free-standing
\makeatletter
\ifx\paragraph\undefined\else
  \let\oldparagraph\paragraph
  \renewcommand{\paragraph}{
    \@ifstar
      \xxxParagraphStar
      \xxxParagraphNoStar
  }
  \newcommand{\xxxParagraphStar}[1]{\oldparagraph*{#1}\mbox{}}
  \newcommand{\xxxParagraphNoStar}[1]{\oldparagraph{#1}\mbox{}}
\fi
\ifx\subparagraph\undefined\else
  \let\oldsubparagraph\subparagraph
  \renewcommand{\subparagraph}{
    \@ifstar
      \xxxSubParagraphStar
      \xxxSubParagraphNoStar
  }
  \newcommand{\xxxSubParagraphStar}[1]{\oldsubparagraph*{#1}\mbox{}}
  \newcommand{\xxxSubParagraphNoStar}[1]{\oldsubparagraph{#1}\mbox{}}
\fi
\makeatother


\usepackage{longtable,booktabs,array}
\newcounter{none} % for unnumbered tables
\usepackage{calc} % for calculating minipage widths
% Correct order of tables after \paragraph or \subparagraph
\usepackage{etoolbox}
\makeatletter
\patchcmd\longtable{\par}{\if@noskipsec\mbox{}\fi\par}{}{}
\makeatother
% Allow footnotes in longtable head/foot
\IfFileExists{footnotehyper.sty}{\usepackage{footnotehyper}}{\usepackage{footnote}}
\makesavenoteenv{longtable}
\usepackage{graphicx}
\makeatletter
\newsavebox\pandoc@box
\newcommand*\pandocbounded[1]{% scales image to fit in text height/width
  \sbox\pandoc@box{#1}%
  \Gscale@div\@tempa{\textheight}{\dimexpr\ht\pandoc@box+\dp\pandoc@box\relax}%
  \Gscale@div\@tempb{\linewidth}{\wd\pandoc@box}%
  \ifdim\@tempb\p@<\@tempa\p@\let\@tempa\@tempb\fi% select the smaller of both
  \ifdim\@tempa\p@<\p@\scalebox{\@tempa}{\usebox\pandoc@box}%
  \else\usebox{\pandoc@box}%
  \fi%
}
% Set default figure placement to htbp
\def\fps@figure{htbp}
\makeatother





\setlength{\emergencystretch}{3em} % prevent overfull lines

\providecommand{\tightlist}{%
  \setlength{\itemsep}{0pt}\setlength{\parskip}{0pt}}



 


\KOMAoption{captions}{tableheading}
\makeatletter
\@ifpackageloaded{caption}{}{\usepackage{caption}}
\AtBeginDocument{%
\ifdefined\contentsname
  \renewcommand*\contentsname{Table of contents}
\else
  \newcommand\contentsname{Table of contents}
\fi
\ifdefined\listfigurename
  \renewcommand*\listfigurename{List of Figures}
\else
  \newcommand\listfigurename{List of Figures}
\fi
\ifdefined\listtablename
  \renewcommand*\listtablename{List of Tables}
\else
  \newcommand\listtablename{List of Tables}
\fi
\ifdefined\figurename
  \renewcommand*\figurename{Figure}
\else
  \newcommand\figurename{Figure}
\fi
\ifdefined\tablename
  \renewcommand*\tablename{Table}
\else
  \newcommand\tablename{Table}
\fi
}
\@ifpackageloaded{float}{}{\usepackage{float}}
\floatstyle{ruled}
\@ifundefined{c@chapter}{\newfloat{codelisting}{h}{lop}}{\newfloat{codelisting}{h}{lop}[chapter]}
\floatname{codelisting}{Listing}
\newcommand*\listoflistings{\listof{codelisting}{List of Listings}}
\makeatother
\makeatletter
\makeatother
\makeatletter
\@ifpackageloaded{caption}{}{\usepackage{caption}}
\@ifpackageloaded{subcaption}{}{\usepackage{subcaption}}
\makeatother
\usepackage{bookmark}
\IfFileExists{xurl.sty}{\usepackage{xurl}}{} % add URL line breaks if available
\urlstyle{same}
\hypersetup{
  pdftitle={Artículo 1: Redes bayesianas multinomiales},
  colorlinks=true,
  linkcolor={blue},
  filecolor={Maroon},
  citecolor={Blue},
  urlcolor={Blue},
  pdfcreator={LaTeX via pandoc}}


% Bibliography setup
\usepackage[
  backend=biber,
  style=apa,
  sorting=nyt
]{biblatex}

\addbibresource{references.bib}



\addbibresource{references.bib}



\title{Uso de Redes Bayesianas multinomiales para la comprensión del uso, seguridad y eficiencia de los medios de transporte en México usando la Encuesta Nacional de Movilidad y Transporte}

\author{
Jessica Milagros Cordero Rivera
\and
Diego Rodrigo González Ramírez
\and
Jose Ángel Lechuga Carrasco
\and
Julio Lizárraga Beltrán
}


\date{30/08/2026}
\begin{document}
\maketitle



\begin{abstract}
Redes Bayesianas Multinomiales (RBM)
\end{abstract}



\section{Introducción}\label{Introducción}
Los medios de transporte son un pilar para el funcionamiento de cualquier sociedad 
moderna, sin embargo, el crecimiento geográfico y económico exponencial que presentan 
los tiempos contemporáneos han provocado complicaciones en estos medios, reduciendo su 
seguridad y eficiencia. Esta situación es especialmente notoria en México, donde la 
urbanización insuficientemente organizada y gran presencia de zonas rurales favorecen 
el caos, haciendo muy complejo el análisis y comprensión de factores de tipo 
demográficos, económicos, geográficos, culturales y sociológicos, que afectan 
negativamente a la eficiencia y/o percepción de estos medios de transporte, 
y por lo tanto, a su uso.

El efecto que tienen estos factores categóricos, cómo la edad, el sexo, la 
escolaridad, el ingreso, el tamaño de la localidad y la interacción 
de estas variables entre si sobre el medio de transporte que optan por 
utilizar los mexicanos, la importancia de la seguridad en el medio 
de transporte y la utilización de esta información para la planificación  
urbana han sido estudiados y confirmados previamente:

Harbering y Schl{\"u}ter estudiaron los factores que afectan la desición en el medio 
de transporte elegido en el Valle de México, mostrando la relevancia de factores 
socioeconómicos y geográficos en las decisiones del medio de tranporte. \cite{Harbering2020}. 
De manera similar, Bautista-Hernández analizó la relación entre las características 
del entorno urbano y la selección de distintos medios de transporte en la Ciudad 
de México y encontraron factores relacionados con las variables socioeconómicas 
y la distancia al centro urbano.
\cite{BautistaHernandez2021}.

Moreno et al. estudiaron patrones de movilidad sostenible en el Área Metropolitana 
de Monterrey y encontraron que los hombres son más propensos a utilizar 
medios de transporte no sostenibles que las mujeres, entre otras cosas.
\cite{Moreno2023}, mientras que Obregón Biosca et al. analizaron la elección 
del medio de transporte entre zonas urbanas y periurbanas y concluyeron que 
los ingresos de los usuarios son una de las variables más significativas en 
la elección del medio de transporte en ambas zonas \cite{ObregonBiosca2018}. 
Estos antecedentes muestran que la movilidad depende de múltiples características que 
interactúan entre sí.

Si bien la literatura previa muestra la clara relevancia de estas 
variables en el medio de transporte mexicano, la mayoría usan análisis
con una variable de salida y varios de entrada, ese método es efectivo, 
pero no encapsula del todo la complejidad de la interacción de estas 
variables entre si, por lo que este trabajo propone un modelo que 
pueda considerar la complejidad de la situación utilizando una Red Bayesiana 
Multinomial (RBM).

Las RBM son modelos probabilísticos que permiten mostrar las relaciones 
de causalidad entre distintas variable categóricas utilizando un 
grafo acíclico dirigido (DAG). La RBM resulta útil para este 
tipo de situaciones, donde el medio de transporte es afectado por 
la interacción de varios factores que normalmente no son independientes.

En este trabajo se utilizan datos de la Encuesta Nacional de Movilidad y Transporte (ENMT) 
para construir una RBM mediante el paquete \texttt{bnlearn} 
del lenguaje de programación R. El modelo busca analizar las relaciones probabilísticas 
entre distintas características de la población mexicana, los medios de transporte 
utilizados y variables relacionadas con la eficiencia y seguridad del transporte. 
Posterior a ello se realizaron cuatro consultas probabilísticas orientadas 
a estudiar cómo determinadas características demográficas, económicas y geográficas 
modifican la probabilidad de utilizar ciertos medios de transporte y la percepción 
de sus condiciones de seguridad y eficiencia.

El objetivo de este análisis Conocer los medios de transporte utilizados en México, 
así como la eficiencia y seguridad del transporte público. Y contar con evidencia 
que ayude a desarrollar políticas sobre planificación urbana.

\section{Metodología}\label{Metodología}

\section{Aplicación}\label{Aplicación}

\subsection{Descripción y preparación de los datos}

Para el análisis se utilizaron variables relacionadas 
con características demográficas y socioeconómicas de los participantes, tamaño de 
la localidad, región, escolaridad, ingreso, condición física, frecuencia de uso de 
distintos medios de transporte y percepción de su eficiencia.

Previo a la construcción de las redes se realizó un proceso de limpieza y 
transformación de los datos. Las variables de interés fueron convertidas a factores y 
algunas variables numéricas fueron recodificadas en categóricas para poder 
realizar el análisis multinomial.Adicionalmente, se creó una variable denominada 
\texttt{preferencia}, construida 
a partir de la frecuencia reportada de uso de diferentes medios de transporte. 
Esta variable permitió resumir la información de varias preguntas de la encuesta.

\subsection{Construcción y comparación de las redes bayesianas}

A partir de las variables seleccionadas se propusieron tres DAG.  
Los DAG propuestos presentan diferentes hipótesis respecto a la dirección y 
dependencia entre las variables. El primer modelo prioriza relaciones entre región, 
tamaño de localidad, ingreso, escolaridad y variables relacionadas con el transporte. 
El segundo incorpora una mayor influencia del tipo de municipio y del ingreso 
individual. El tercer modelo propone una estructura más tranquila, en la que 
la región, escolaridad, tamaño de localidad y condición física funcionan como 
variables relevantes para explicar los patrones de movilidad observados.

\begin{figure}[H]
    \centering

    \begin{subfigure}[b]{0.31\textwidth}
        \centering
        \includegraphics[width=\textwidth]{DAG1.jpeg}
        \caption{DAG 1}
        \label{fig:dag1}
    \end{subfigure}
    \hfill
    \begin{subfigure}[b]{0.31\textwidth}
        \centering
        \includegraphics[width=\textwidth]{DAG2.jpeg}
        \caption{DAG 2}
        \label{fig:dag2}
    \end{subfigure}
    \hfill
    \begin{subfigure}[b]{0.31\textwidth}
        \centering
        \includegraphics[width=\textwidth]{DAG3.jpeg}
        \caption{DAG 3}
        \label{fig:dag3}
    \end{subfigure}

    \caption{Grafos acíclicos dirigidos propuestos para representar las relaciones 
    de dependencia entre las variables de interés.}
    \label{fig:dags}
\end{figure}

\subsection{Validación de los modelos}

Para determinar cuál de las tres estructuras propuestas representaba mejor los 
datos, los modelos fueron comparados mediante los criterios AIC y BIC utilizando 
las funciones implementadas en el paquete \texttt{bnlearn}.

\begin{table}[H]
\centering
\caption{Comparación de los DAG propuestos mediante BIC y AIC.}
\label{tab:aicbic}
\begin{tabular}{lcc}
\toprule
Modelo & BIC & AIC \\
\midrule
DAG 1 & [-5547.98] & [-4678.218] \\
DAG 2 & [-13916.51] & [-7193.286] \\
DAG 3 & [-6233.563] & [-4816.989] \\
\bottomrule
\end{tabular}
\end{table}

Con esta información podemos concluir que el modelo de la DAG 1 es el que mejor
ajusta las relaciones de dependencia de nuestras variables categóricas

\subsection{Creación de una DAG mediante "Hill climbing"}

Adicional a los 3 DAG realizados, se utilizó el algoritmo de \texttt{hill climbing} 
dentro de \texttt{bnlearn} para hacer el mejor modelo estadístico.

\begin{figure}[h]
    \centering
    \includegraphics[width=1\linewidth]{BESTDAG.jpeg}
    \caption{DAG hecha con el algoritmo "Hill climbing"}
    \label{bestdag}
\end{figure}

Sin embargo, este modelo fue descartado rápidamente debido a la poca dependencia que hay 
entre las variables de interés.

\subsection{Consultas probabilísticas}

\subsubsection{Percepción de rapidez según el tamaño de la localidad}

La primera consulta evaluó si la percepción de rapidez del transporte presentaba 
diferencias entre habitantes de localidades clasificadas como urbanas y aquellos 
pertenecientes al resto de las localidades consideradas en el análisis. Para ello 
se construyó un indicador de percepción de rapidez a partir de las evaluaciones 
de diferentes medios de transporte.

La probabilidad estimada de reportar una buena percepción de rapidez fue de 
$0.6682$ para el grupo clasificado como urbano y de $0.7271$ para el grupo 
clasificado como rural. Los resultados muestran, por lo tanto, una mayor 
probabilidad de evaluar favorablemente la rapidez del transporte dentro del 
segundo grupo. Este resultado sugiere que el contexto territorial puede estar asociado con la 
percepción del funcionamiento de los medios de transporte.


\subsubsection{Escolaridad y percepción de eficiencia}

La segunda consulta analizó la relación entre escolaridad y percepción de eficiencia 
del transporte público. La escolaridad fue agrupada en dos categorías, básica y 
alta, mientras que la percepción fue clasificada como eficiente o ineficiente.

La probabilidad de considerar eficiente al transporte fue de $0.6440$ entre las 
personas clasificadas con escolaridad básica y de $0.5897$ entre aquellas con 
escolaridad alta. Contrario a la hipótesis de que una mayor escolaridad estaría 
asociada con una mejor percepción de eficiencia, los datos mostraron una 
probabilidad mayor de evaluación favorable dentro del grupo con escolaridad básica.

\subsubsection{Uso del transporte y percepción de eficiencia}

La tercera consulta evaluó si la percepción de eficiencia cambia entre las personas 
que utilizan habitualmente los medios de transporte considerados y aquellas que 
no presentan este patrón de uso.

Entre las personas clasificadas como usuarias, la probabilidad de considerar 
eficiente al transporte fue de $0.6612$. En contraste, para las personas 
clasificadas como no usuarias la probabilidad obtenida fue de $0.4560$.

Este resultado muestra una diferencia muy 
marcada entre los grupos analizados. Por lo tanto se puede concluir que 
la experiencia de uso se encuentra asociada con la percepción de eficiencia 
del transporte dentro de la muestra. 


\subsubsection{Condición física y uso cotidiano del automóvil}

La cuarta consulta estudió la relación entre condición física y frecuencia de uso 
del automóvil dentro del subconjunto correspondiente a las localidades de mayor 
tamaño.

La probabilidad estimada de utilizar el automóvil cotidianamente fue de $0.5972$ 
para las personas clasificadas con incapacidad y de $0.7326$ para aquellas sin 
incapacidad. Aunque los valores obtenidos muestran diferencias descriptivas entre 
ambos grupos, el tamaño de muestra disponible para esta comparación fue insuficiente.


\begin{table}[H]
\centering
\caption{Resultados de las consultas probabilísticas.}
\label{tab:queries}
\begin{tabular}{p{6cm}lc}
\toprule
Consulta & Evidencia & Probabilidad \\
\midrule

Buena percepción de rapidez 
& Localidad urbana & 0.6682 \\

& Localidad rural & 0.7271 \\
\midrule

Percepción de transporte eficiente 
& Escolaridad básica & 0.6440 \\

& Escolaridad alta & 0.5897 \\
\midrule

Percepción de transporte eficiente 
& Usuario & 0.6612 \\

& No usuario & 0.4560 \\
\midrule

Uso cotidiano del automóvil 
& Con incapacidad & 0.5972 \\

& Sin incapacidad & 0.7326 \\

\bottomrule
\end{tabular}
\end{table}

Los resultados obtenidos muestran que las características territoriales, educativas 
y de uso del transporte se encuentran asociadas con cambios en las probabilidades 
de percepción y utilización de los distintos medios considerados. En particular, 
la experiencia de uso presentó una relación importante con la percepción de 
eficiencia, mientras que la escolaridad mostró un comportamiento diferente al 
planteado inicialmente.

Las redes bayesianas permitieron representar estas relaciones dentro de un modelo 
probabilístico y responder preguntas condicionadas a diferentes tipos de evidencia. 
A diferencia de una comparación exclusivamente descriptiva, este enfoque permite 
actualizar la probabilidad de una variable de interés a partir de las características 
observadas en otros nodos de la red.

No obstante, los resultados deben interpretarse como asociaciones probabilísticas 
y no como efectos causales. Las estructuras utilizadas fueron construidas a partir 
de hipótesis de dependencia y datos observacionales, por lo que la presencia de un 
arco entre dos variables no demuestra que una de ellas cause directamente a la otra.

\section{Conclusiones}\label{Conclusiones}

\printbibliography


\end{document}